\section{Experimental Setup}
\label{sec:setup}

This section defines the elements shared by all experiments: the datasets, the
comparability-preserving test-set filtering, the model and training protocol, and the
evaluation metrics. Experiment-specific choices are described in the respective
sections (\Cref{sec:gap,sec:weather-aug,sec:pseudo-labels,sec:style-transfer}).

\subsection{Datasets}
\label{subsec:setup-data}

\paragraph{Synthetic training sources (LARD-V2).}
Training uses the five synthetic sources of LARD-V2 --- \texttt{flsim} (Flight
Simulator), \texttt{arcgis} (ArcGIS), \texttt{bingmaps} (Bing Maps), \texttt{ges}
(Google Earth Studio) and \texttt{xplane} (X-Plane) --- generated for approximately
\textbf{250 airports worldwide}. The full synthetic training pool comprises
\textbf{66\,565 images}. To make the sources directly comparable, the images are
grouped so that every runway viewpoint is represented by the same number of images in
each source; this grouping reduces the pool to \textbf{7\,715 images per source}
(\textbf{38\,575} in total). A key consequence of this design is that the individual
sources share \texttt{location\_id}s: image $i$ in each source depicts the \emph{same}
runway view rendered by a different pipeline, which turns the cross-source comparison
of \Cref{sec:gap} into a controlled experiment in which the rendering source is the
only variable.

\paragraph{Training budget and the combined source.}
From this pool, each training run fine-tunes on a fixed budget of \textbf{5000
images}. The \texttt{combined} source (also referred to as \texttt{all}) mixes all
five pipelines in equal proportion at matched locations; to respect the 5000-image
budget it therefore covers roughly \emph{one fifth} as many distinct locations as an
individual source (each location contributing five renderings instead of one). This
sets up an explicit trade-off between \emph{rendering-source diversity} and
\emph{runway-viewpoint coverage} that recurs throughout the thesis. An extended
variant that relaxes the budget, \texttt{all\_extended}, is included where noted.

\paragraph{Real test set (LARD-V1) and weather splits.}
Real-world performance is evaluated on a held-out set of \textbf{1880 real images}
from LARD-V1, composed of \textbf{1557 nominal} images and \textbf{323 edge cases}.
The edge cases, together with the nominal split, are stratified into five
weather/lighting scenarios: \emph{nominal}, \emph{glare}, \emph{rain}, \emph{fog} and
\emph{night}. Because the nominal scenario accounts for almost three quarters of the
set, the \emph{overall} real metric is nominal-dominated; per-scenario numbers are
therefore used whenever adverse conditions are discussed. To enable a fair
comparison against the synthetic domain, an equally sized sample is drawn from the
synthetic sources (from images not used for training). \Cref{tab:weather-splits}
summarizes the composition.

\begin{table}[th]
    \caption{Real test set (LARD-V1), composition and weather stratification. The
    nominal split dominates, which is why per-scenario metrics are reported for
    adverse conditions.}
    \label{tab:weather-splits}
    \centering
    \begin{NiceTabular}{lr}
        \CodeBefore
        \rowcolors{2}{gray!10}{white}
        \Body
        \toprule
        \textbf{Split} & \textbf{\# Images} \\
        \midrule
        Nominal        & 1557 \\
        Edge cases (glare + rain + fog + night) & 323 \\
        \midrule
        \textbf{Total} & 1880 \\
        \bottomrule
    \end{NiceTabular}
\end{table}

\paragraph{Unlabeled real data and the custom test set (self-training).}
The self-training experiments (\Cref{sec:pseudo-labels}) additionally use
\emph{unlabeled} real landing sequences collected from public YouTube channels,
captured at airports and channels that are \emph{disjoint} from the LARD-V1 test set.
The sequences span the same adverse-weather scenarios as the test splits (nominal,
glare, rain, fog) and are decomposed into approximately \textbf{1800 frames}. A
subset of these sequences forms a \emph{custom} landing-sequence test set (with
glare, fog and rain splits) used to complement the LARD-V1 evaluation.

\subsection{Test-Set Filtering for Comparability}
\label{subsec:setup-filtering}
A naive comparison between the synthetic and real test sets is confounded by
differences in the \emph{apparent runway size} and \emph{orientation} of their
runways. This thesis therefore filters the test sets along two axes before comparing
them:
\begin{itemize}
    \item \textbf{Runway size.} The apparent runway size is measured as the
    $\log_{10}$ of the fraction of image pixels occupied by the runway, and the test
    sets are filtered to a common size distribution.
    \item \textbf{Vertical angle.} Runways in the synthetic test sets frequently
    appear at oblique orientations, whereas real approaches are mostly aligned
    vertically. The \emph{vertical angle} is defined as the absolute value of the
    angle of the runway with respect to the image $y$-axis, taking values in
    $[0^\circ, 90^\circ]$; it measures the runway's ``skewness'' independently of
    direction. Matching this distribution across test sets removes an orientation
    bias that otherwise flatters the real set.
\end{itemize}
The effect of adding the vertical-angle criterion on top of the size criterion is
analyzed in \Cref{sec:gap}.

\subsection{Model}
\label{subsec:setup-model}
All experiments fine-tune a pretrained \textbf{YOLOv8-pose} model with four keypoints
corresponding to the runway corners. Fine-tuning from a pretrained checkpoint (rather
than training from scratch) keeps the setup lightweight and reproducible and ensures
that differences in results are attributable to the training data and adaptation
technique rather than to backbone initialization. This choice is further motivated by
the LARD-V2 benchmark, on which the YOLO family performs strongly
(\Cref{subsec:rw-runway}).
\todo[inline]{State the exact variant (n/s/m/l), input resolution, optimizer,
learning-rate schedule, batch size and hardware.}

\subsection{Training Protocol}
\label{subsec:setup-protocol}
Two training budgets are used, depending on the experiment:
\begin{itemize}
    \item \textbf{Full budget (default).} 200 epochs with early-stopping
    \emph{patience} of 25 epochs. Used for the Sim2Real gap characterization
    (\Cref{sec:gap}), self-training (\Cref{sec:pseudo-labels}) and style transfer
    (\Cref{sec:style-transfer}).
    \item \textbf{Reduced budget (weather-augmentation sweep).} 20 epochs, used only
    for the weather-augmentation study (\Cref{sec:weather-aug}) to keep the large
    augmentation$\times$ratio sweep tractable.
\end{itemize}

\paragraph{Sampling variance and repeated runs.}
Two sources of randomness strongly affect the mAP of a single run: (i) the specific
images sampled into the fixed 5000-image budget (especially for the \texttt{combined}
source, which re-samples locations), and (ii) the specific images selected to receive
an augmentation. To control for this, each configuration in \Cref{sec:weather-aug} is
trained with \textbf{4 independent runs}. As detailed there, the \texttt{combined}
training set is additionally \emph{fixed} (sampled four times, best retained) so that
the reported variance is attributable to the augmentation alone.

\subsection{Evaluation Metrics}
\label{subsec:setup-metrics}
The \textbf{primary metric} is keypoint \ac{mAP} averaged over the ten \ac{OKS}/\ac{IoU}
thresholds from $0.50$ to $0.95$ in steps of $0.05$, denoted \textbf{mAP@[.50:.95]}
(equivalently mAP50--95). It is the strictest single-number summary and rewards
accurate, well-localized runway quadrilaterals rather than mere presence.
mAP@$0.50$ is reported alongside it as a looser reference. Where a single localization
figure is more informative --- for example the runway-size analyses of
\Cref{sec:gap} --- the \textbf{\ac{IoU}} between the predicted and ground-truth runway
region is reported directly. Pixel-level precision/recall of the runway corners and
\textbf{inference speed} (throughput, frames per second on the reference hardware) are
reported as secondary, deployment-oriented metrics, in line with the edge-deployable
target application.

All metrics are reported \emph{overall} (entire real test set) and \emph{per weather
scenario}; because the overall figure is nominal-dominated, per-scenario numbers are
used when reasoning about adverse conditions.

\paragraph{What counts as a meaningful improvement.}
Given the four-run variance-control protocol, an adaptation strategy is considered to
\emph{meaningfully} improve over the synthetic-only baseline on a given split when its
mean mAP@[.50:.95] gain exceeds the run-to-run spread of the baseline. As a practical
guide, single-point gains ($\sim$1\,mAP) are treated as marginal, whereas gains of
several mAP points on the adverse-weather splits --- where the gap is largest --- are
treated as substantive.
\todo[inline]{Fix the OKS keypoint scales for the 4-corner runway pose, the detection
confidence/IoU thresholds, and the exact spread value used for the ``meaningful
improvement'' criterion once baseline variance is measured.}
